\chapter{计算机的层次}
\label{ch:what-a-computer-does}
\chapterlead{本章说明一台计算机从通电到出现光标之间发生了什么。这些层次解释了后续所有命令的形态。}

\section{起步：查看本机配置}

无论在 Windows、macOS 还是已经安装的 Linux 上，先取得四个数字：处理器核心数、内存容量、磁盘剩余空间、系统与内核版本。后续排查性能与容量问题时，需要这些基线。

\begin{itemize}
  \item Windows：任务管理器 \textrightarrow\ 性能，可看到 CPU、内存、磁盘的实时曲线；\texttt{Win+R} 输入 \texttt{msinfo32} 查看系统摘要。
  \item macOS：活动监视器，标签页分别对应 CPU、内存、磁盘、网络；“关于本机”显示芯片与内存。
  \item Linux：运行下列四条命令。
\end{itemize}

\begin{shellsession}
$ uname -a
Linux nixos 7.3.4 #1-NixOS SMP PREEMPT_DYNAMIC x86_64 GNU/Linux
$ nproc
8
$ free -h
               total        used        free
Mem:            31Gi        9.8Gi        12Gi
$ df -h /
Filesystem      Size  Used Avail Use% Mounted on
/dev/nvme0n1p2  476G  182G  270G  41% /
\end{shellsession}

\texttt{uname -a} 给出内核名、内核版本、构建信息与体系结构；\texttt{nproc} 给出可用核心数；\texttt{free -h} 以 GB 为单位显示内存；\texttt{df -h /} 显示根分区的容量与占用比例。这四个数值记在笔记首页。机器变慢或磁盘告警时，先用它们与当前值比较。

\section{硬件：执行指令的四个部件}

机箱内的运算部件有四类。

\begin{itemize}
  \item \textbf{处理器（CPU）}：执行指令。每秒可执行数十亿条，但每条指令的操作极简单：相加、比较、搬移。
  \item \textbf{内存（RAM）}：存放运行中的程序与数据。断电后内容丢失。
  \item \textbf{磁盘（SSD 或硬盘）}：长期存放数据，断电后内容保留。
  \item \textbf{总线与设备}：连接上述部件与屏幕、键盘、网卡。
\end{itemize}

一段程序由两部分组成：指令序列，以及指令处理的数据。屏幕上的每个像素都是内存中的一个数值；文字、图片、视频都是这些数值按约定排列的结果。硬件不识别字符，只搬运与计算数值。

\begin{figure}
\centering
\figfit{%
\begin{tikzpicture}[font=\small, node distance=8mm]
  \node[draw=rule, fill=white, minimum width=132mm, minimum height=11mm, align=center]
    (apps) {应用程序\\[-1pt]{\footnotesize\color{muted} 浏览器 \textperiodcentered\ 编辑器 \textperiodcentered\ 终端中的命令}};
  \node[draw=ochre!55, fill=ochre!8, minimum width=132mm, minimum height=8mm, below=of apps, align=center]
    (sysc) {系统调用接口};
  \node[draw=teal!55, fill=teal!9, minimum width=132mm, minimum height=14mm, below=of sysc, align=center]
    (kern) {内核\\[-1pt]{\footnotesize\color{muted} 进程调度 \textperiodcentered\ 内存管理 \textperiodcentered\ 文件系统 \textperiodcentered\ 设备驱动}};
  \node[draw=rule, fill=white, minimum width=132mm, minimum height=9mm, below=of kern, align=center]
    (hw) {硬件\\[-1pt]{\footnotesize\color{muted} 处理器 \textperiodcentered\ 内存 \textperiodcentered\ 磁盘 \textperiodcentered\ 网卡}};
  \draw[-{Latex[length=1.5mm]}, ochre, thick] ([xshift=-8mm]apps.south) -- ([xshift=-8mm]sysc.north);
  \draw[-{Latex[length=1.5mm]}, muted] ([xshift=8mm]sysc.north) -- ([xshift=8mm]apps.south);
  \draw[-{Latex[length=1.5mm]}, teal, thick] ([xshift=-8mm]sysc.south) -- ([xshift=-8mm]kern.north);
  \draw[-{Latex[length=1.5mm]}, teal, thick] ([xshift=-8mm]kern.south) -- ([xshift=-8mm]hw.north);
\end{tikzpicture}
}
\caption{机器的四层结构。程序不直接访问硬件，请求逐层向下，结果逐层返回。}
\end{figure}

\section{操作系统：资源竞争的仲裁者}

一台机器上通常同时运行数百个程序。浏览器等待网络，播放器向声卡写数据，终端中的命令读写文件，系统自身检查温度。这些程序竞争处理器时间、内存空间与磁盘访问。

如果程序可以直接操作硬件，一个地址计算错误就会覆盖其他程序的数据，一个死循环会使整台机器失去响应。操作系统的职责是在这些竞争中维持秩序，具体分为四项：进程调度、内存分配、文件组织、设备访问。这四项构成第一部分的主线。

处理器在同一时刻只执行一个程序的一段指令。“同时运行数百个程序”的实现方式是在极短的时间片内反复切换。内存管理采用类似方法：每个程序认为自己独占一块连续地址空间，实际访问每个地址都要经过操作系统翻译，越界访问被拒绝。

\section{内核与用户态}

操作系统最核心的部分是内核（kernel）。内核运行在处理器的高特权模式，可以直接访问内存与硬件。普通程序运行在低特权模式，不能直接写磁盘、配置网卡或终止其他进程。

程序需要内核服务时，通过系统调用（system call）提出请求。执行 \texttt{rm} 删除文件时，\texttt{rm} 本身不擦除磁盘数据；它请求内核执行删除，内核检查权限后再决定是否执行。

这一设计产生两个结果。第一，权限由内核强制执行，不是界面提示。普通用户读不到他人的私有文件，原因是内核拒绝了请求。第二，最高权限账户（root）不受多数限制，可以删除系统文件、修改内核参数、读取任何用户的数据。长期以 root 身份登录会放大误操作的后果：一条错误的删除命令可能使系统无法启动。本书的操作默认在普通用户下执行，需要时临时提权。

\begin{warnbox}[不要长期以 root 登录]
root 身份下没有任何机制阻止破坏性操作。初学者最常见的事故是在 root shell 中执行了本应在特定目录内执行的删除命令。第\ref{ch:permissions}章说明 \texttt{sudo} 的正确用法。
\end{warnbox}

\section{Linux 的构成}

Linux 严格意义上仅指一个内核。1991 年，林纳斯·托瓦兹发布了一个可在个人计算机上运行的内核。仅有内核无法使用：列目录、编辑文件、安装软件都需要相应的程序。

这些程序大多来自 GNU 项目，因此也有 GNU/Linux 的称法。日常语境中的 Linux 指内核、GNU 工具与桌面软件的整体。

发行版（distribution）把内核、工具、包管理器、默认配置与社区支持打包在一起。Ubuntu、Debian、Fedora、Arch、NixOS 都是发行版，共用同一内核，差异在于随附工具、软件安装方式、默认配置与支持渠道。

\begin{notebox}[图形界面不属于内核]
图形界面是独立的一层软件：显示协议（X11 或 Wayland），以及负责绘制窗口、处理输入的程序，即桌面环境或合成器，例如 GNOME、KDE Plasma、Hyprland、Niri。同一套 Linux 可以搭配其中任意一个，命令行部分不受影响。
\end{notebox}

\section{发行版的选择}

可选发行版有数百个，对初学者有实际差别的只有少数几个。下列对比只列影响学习成本的部分。

\ifbknarrow
% 五列表在窄页上挤成一团，改用列表，信息不减
\begin{itemize}
  \item \textbf{Ubuntu LTS}：两年一个长期支持版，apt。教程最多，代价是默认配置多、跨版本升级易出问题。
  \item \textbf{Debian}：约两年，更新保守，apt。适合长期稳定的服务器，软件版本偏旧。
  \item \textbf{Fedora}：约半年，dnf。新硬件支持好，生命周期短、升级频繁。
  \item \textbf{Arch}：滚动更新，pacman。文档质量高，冲突需自行处理。
  \item \textbf{NixOS}：半年一个稳定版，nix。可回滚、可复现，学习曲线陡、教程不通用。
\end{itemize}
\else
\begin{center}
\centering
\small
\begin{booktable}{\textwidth}{L{0.783} L{1.160} L{0.720} L{1.493} J{0.844}}
\caption{五个常见发行版的取舍}\\
\toprule
发行版 & 更新节奏 & 包管理 & 适用对象 & 主要代价 \\
\midrule
\endfirsthead
\toprule
发行版 & 更新节奏 & 包管理 & 适用对象 & 主要代价 \\
\midrule
\endhead
Ubuntu LTS & 两年一个长期支持版 & apt & 需要最多教程与第三方软件 & 默认配置多，跨版本升级易出问题 \\
Debian & 约两年，更新保守 & apt & 需要长期稳定的服务器 & 软件版本偏旧 \\
Fedora & 约半年 & dnf & 使用新硬件、需要较新软件 & 生命周期短，升级频繁 \\
Arch & 滚动更新 & pacman & 愿意阅读文档并自行维护 & 升级冲突需自行处理 \\
NixOS & 半年一个稳定版 & nix & 需要理解系统构成与回滚能力 & 学习曲线陡，教程不通用 \\
\bottomrule
\end{booktable}
\end{center}
\fi

Ubuntu 的第三方教程数量最多，代价是预装与预配置项也最多。Debian 更新保守，在服务器上常见。Fedora 对新硬件支持较好。Arch 的文档（Arch Wiki）质量高，前提是使用者愿意自行维护系统。

\section{本书的选择}

传统发行版的状态是累积形成的。按教程安装软件、修改配置文件、添加第三方软件源、升级系统，半年之后没有人能完整说明这台机器为何可运行。在新机器上复现同样的环境只能凭记忆重做；回到三个月前的状态则无法实现。

NixOS 把系统状态写成配置。安装的软件、启用的服务、用户、网络、桌面都记录在配置文件中。执行一次重建，系统按配置生成新状态，切换后生效。每次重建留下一个代（generation），出现问题时用一条命令回到上一代。

支撑这一机制的是只读软件仓库 \texttt{/nix/store}。每个软件包存放在以哈希命名的目录中，安装不覆盖旧版本，卸载不影响正在运行的程序。

代价有四项：需要学习一小部分 Nix 语言；多数网上教程不能照搬，教程中“编辑 \texttt{/etc} 下某个文件”的步骤在此处应改为修改配置并重建；磁盘占用高于传统发行版，同一软件的多个版本并存，需要定期清理；未被 nixpkgs 收录的软件需要自行编写打包表达式。

对本书的目标读者，代价与收益的交换如下：需要学习的内容是系统的构成，而 NixOS 把这部分以配置的形式显式呈现。回滚机制也降低了尝试的成本。若当前目标仅是在一周内让某个软件运行，Ubuntu 更省时。两种选择并不冲突。

本章只给出结论。第\ref{ch:why-nixos}章展开这两种系统状态的来源、可复现的范围、六项代价，以及这些性质在 AI Agent 参与系统维护之后的位置。

\section{本书的路径}

读者手上的机器不同：可能有一台可整机安装的旧电脑，也可能只有一台正在使用的笔记本，或者租用了一台服务器。三条路线分别成章，可按条件选择。本书推荐实体机：它给出完整的系统，显卡、无线、休眠这些内容只有在真实硬件上才会遇到。没有可整机安装的机器时，虚拟机的流程与之一致，见第\ref{ch:install-vm}章。

学习顺序不随路线变化：先熟练命令行，再理解 NixOS 的配置方式，最后进入服务器与 AI 工具。跳过前两步直接部署服务器，通常会在第一次故障时停滞。阅读指南中有详细路线表与排版约定。

\section{练习}

\begin{enumerate}
  \item 查清本机的处理器型号、内存容量与磁盘剩余空间，记录在笔记首页。
  \item 打开任务管理器或活动监视器，统计当前进程数，选出三个常驻后台的进程，说明它们的用途。
  \item 用一个段落说明：普通程序为什么不能直接格式化磁盘。
  \item 用一个段落说明：与 Ubuntu 相比，NixOS 在“系统可复现”这一项上的差异体现在哪里。
\end{enumerate}

\begin{checkpoint}
本章结束时，应能：
\begin{itemize}
  \item 说明处理器、内存、磁盘各自承担的功能；
  \item 说明操作系统为何需要管理进程、内存、文件与设备；
  \item 区分内核与普通程序的权限，说明 root 的风险来源；
  \item 描述 Linux 内核、GNU 工具与发行版的关系；
  \item 说明本书选择 NixOS 的收益与四项代价。
\end{itemize}
\end{checkpoint}
